\section{Self Knowledge}

\begin{quotex}
I am worm, I am slave, I am king, I am God. \flright{\textsc{Gabriel Derjavine}}

The ``to be silent" therefore precedes the ``to know", the ``to will" and the ``to dare". This is why the Pythagorean school
prescribed five years silence to beginners or ``hearers". One dared to speak there only when one ``knew" and ``was able
to", after having mastered the art of being silent—that is to say, the art of concentration. The
prerogative ``to speak" belonged to those who no longer spoke automatically, driven by the game of the intellect and
imagination, but who were able to suppress it owing to the practice of interior and exterior silence, and who knew what
they were saying—again thanks to the same practice. The silentium practised by Trappist monks and
prescribed for the time of ``retreat", generally to all those there who are taking part, is only the application of the
same true law: ``Yoga is the suppression of the oscillations of the mental substance" or ``concentration is the willed
silence of the automatism of the intellect and imagination". \flright{\textsc{Valentin Tomberg}}

\end{quotex}
The theme of this past Sunday's sermon was self-knowledge. It began with the notion that we are to
know, love, and serve God. The priest pointed out that we cannot serve what we don't love.
Moreover, we cannot love what we don't know (except perhaps in an abstract way). These are the
three forms of a Christian yoga: jnani yoga (to know God), bhakti yoga (to love God), and karma yoga (to serve God).
They are also analogous to the three rules of Hermetism:

\begin{itemize}
\item To know 
\item To will 
\item To dare 
\end{itemize}
Although ``to be silent" is usually the fourth rule, Tomberg puts it first. Interior silence is the ``suppression of the
oscillations of the mental substance". As we learned from Plotinus, the Intellective Act cannot be known unless the
soul is able to reflect it; that is, it must be free of arbitrary oscillations.

\paragraph{Knowledge of the Real I}
Hence, to know God, one must first know oneself in order to find and eliminate the sources of these perturbations of the
soul. What separates Hermetism from philosophy is that the former depends on certain exercises, not being limited to
the discursive thought of the latter. For example, the traditional teaching is that the human being has, besides a
physical body, three levels of the soul: the intellectual soul, the sensitive soul, the nutritive soul (although the
names may be translated differently). In our work, we may use different designations for the same thing, as seen in Table \ref{tab:LayersSoul}:

\begin{table}[h]
\begin{tabular}{ccc}\toprule
Intellectual soul &
Mental body &
Intellectual Centre\\
Sensitive or animal soul &
Astral body &
Emotional Centre\\
Nutritive or vegetative soul &
Etheric body &
Motor Centre\\\bottomrule
\end{tabular}
\caption{Layers of the Soul}
\label{tab:LayersSoul}
\end{table}

These, then, are known, not through philosophy, science or reasoning, but rather through direct observation. To know
oneself, one has to understand how each of these centres operates, and how they interoperate, in
one's own life. However, the soul forms a unity, at least virtually, so there is no sharp
separation. A deeper observation shows how parts of each of the souls interpenetrate each other. Thus, there is
intelligence in the sensitive and nutritive souls, as well as automatic and emotional parts of the intellectual soul.

The first stage of self-knowledge is to recognize that our soul life is fragmented into a multiplicity of
I's, often in conflict with each other. \textbf{Thomas Aquinas} says that ``everything is knowable
so far as it is in act, and not, so far as it is in potentiality [or virtual]" (Question 87). Hence, the soul cannot be
known as long as it is only a virtual unity. The Real I, which exists only in potential, must be actualised. The
actualisation of this Real I is an act of concentration that unites all the apparent disparate parts of the soul.\footnote{There are higher centres than those enumerated here, but they are not the topic of this post.}

\paragraph{Levels of Being}
One then sees that there are four levels of being, corresponding to differing understanding of the I.

\begin{table}[h]
\begin{tabular}{cc}\toprule
Absolute Consciousness &
The Universal I\\
Consciousness of the Real I &
The real individual I\\
Waking Consciousness &
The personal I (actually, a multitude)\\
Subconsciousness &
The physical I of the body\\\bottomrule
\end{tabular}
\caption{Degrees of Consciousness}
\label{tab:DegreesConsciousness}
\end{table}

The \textbf{subconscious} is active, for example, during sleep or in the autonomous functions of the organism. It also
refers to those inner acts that affect our soul life, even when we are not conscious of them.

\textbf{Waking consciousness} is our normal state throughout the day, as it is what we are given as fallen human beings.
However, self-observation shows that this state is really not fully conscious. We can see that it seldom rises above
daydreams the operate just under awareness. It is subjective, leading to features like partisanship, querulousness,
confabulation, fanaticism. The deleterious effects on social life are the results of a population stuck in this state.

\textbf{Self consciousness} is not given to us naturally, except perhaps in unusual conditions like boundary situations
or in activities requiring extreme concentration. This is the state of objective consciousness, that of the just man.

\textbf{Absolute consciousness} is the consciousness of the Absolute. This is how we know God. The Beatific Vision is to
know God in His essence.

This state is related to the notion of \emph{theosis}. We can start with the attributes of God, the \textbf{unmoved
mover}, for example. If the actualisation of the I is deep enough, one can get a sense of what it means to be an
unmoved mover: the I observes, but is not disturbed by external events.



\flrightit{Posted on 2018-10-10 by Cologero }
